\section{Evaluation}\label{sec:evaluation}

The evaluation separates four questions: whether the implementations agree on
the checked core, whether one frozen implementation preserves application
outputs, whether source-derived byte envelopes pass the admission certificate,
and which native observations survive a single-batch serving prototype.
All inference is local. The full application study and the three additional
experiments below have separate protocols and denominators.

\subsection{Core checks and implementation identity}

All 47 configured tests pass, including source binding, immutable resumption,
model-indexed reservations, unsupported stop declarations, complete token-piece
decoding and pre-dispatch byte guards. Actual pinned GGUFs enable the tokenizer
and template tests; model-free runs report those cases as skips.
Generated differential testing retains 1,000 authored programs, 5,000
original/candidate Lean checks and 16,000 trace/store comparisons without a
counterexample. A separate operational comparison covers 150 programs in four
contexts under two finite kernels and five variants, plus a targeted template
program. Python, an independent interpreter and Lean's actual \code{exec} agree
on 6,320 executions and 604 validator checks. These audit worlds use character
counts. They test implementation agreement without verifying the frontend or
native tokenizer. Retained fault injection rejects all 103 mutants that leak
under its finite oracle, among 210 mutants over 12 operators.

Before application inference, we froze production code at \code{41a7d81} and ran
from a clean detached checkout. All nine main workflows and three Qwen variants
use that identity. Manifests bind the complete implementation, source, data,
selected IDs, model contracts and native-library hashes. The study uses Python
3.12.14, llama-cpp-python 0.3.36, native build 11490 (\code{9c2e0e491}), four CPU
threads, a 4,096-token context, batch size 128 and greedy decoding. GPU execution
uses Vulkan on an NVIDIA RTX 5050 Laptop GPU, with 99 requested offload layers.
Seeds are fixed by dataset index and useful-call site. We report logical usage;
this single-machine protocol supports no throughput or timing-protection claim.

\subsection{Full homogeneous application study}\label{sec:final-validation}

The main study contains 1,172 paired inputs over routing, chaining,
evaluator--optimizer, orchestrator--worker, switch, email, calculator, corrective
RAG and adaptive RAG. It uses Llama 3.2 1B, with Llama 3.2 3B for the adaptive
workflow's larger-model path. Replication contains 512 Qwen 2.5 1.5B pairs over
switch, email and corrective RAG. The bounded ports retain their declared loop
limits, retrieval corpora and primary grading prompts. Failed or mismatched
attempts cannot be removed from the completion denominator.

All \roundTwoPairs{} pairs complete with equal saved fields, one shaped trace
per workflow and no recorded truncation. The independent record audit checks
3,368 execution records and \roundTwoCalls{} useful and cover model calls.
It re-tokenizes every complete request, reconstructs visible prefixes and
reconciles per-model ledgers. It also matches private local eval/yield records
to public events, including duplicate multiplicities. A deliberately changed input count is
detected. These records support the declared client accounting contract; they
are not provider attestations or a proof of every intermediate store value.

\begin{table}[t]
\centering\footnotesize
\setlength{\tabcolsep}{4pt}
\caption{One frozen implementation. $K$ is the observed decision-class count;
original balanced accuracy (BA) and task accuracy are percentages. Tr. and MI
are distinct shaped traces and empirical trace/decision mutual information in
bits. Cost is shaped/original weighted logical usage. Q denotes Qwen replication;
a dash omits an unavailable task scorer or an uninformative single-class attack.}
\label{tab:round2-workflows}
\begin{tabular}{lrrrrrrr}\toprule
Workflow & Pairs & $K$ & Orig BA & Tr. & MI & Cost & Task acc.\\\midrule
\tabinput{generated/round2-workflows.tex}
\bottomrule\end{tabular}
\end{table}

Attack features retain ordered model/tool identities, counts, event count and
model-call count. A class-balanced random forest uses five-fold cross-validation,
stratified when class counts permit. These are descriptive classifiers on the
authored corpus. Original switch prediction is near chance, and the primary
Qwen corrective-RAG port has one observed decision class. Neither supplies a
general leakage result. Every shaped feature vector is constant within its
workflow, giving empirical trace/decision MI zero. Fold-dependent prediction
scores can still differ from $1/K$ because the trained constant prediction
depends on each fold's labels; this does not reveal information in a constant
trace.

Output preservation retains task errors. Main-study routing, switch, calculator,
corrective RAG and adaptive RAG achieve 65.3, 37.6, 2.0, 34.0 and 53.3 percent,
respectively, using the inherited scorers. Calculator remains a control-flow
stress test. Preserving these bounded ports does not recover an upstream
application's accuracy or establish independent response quality.

\subsection{Cost and observer alternatives}\label{sec:costs}

For input weight one and output weight four, main-study median overhead is
\roundTwoMedian{} times original usage, with range
\roundTwoMin{}--\roundTwoMax{}. These weights describe logical usage, not a
provider invoice. Static recompilation on these records gives zero median
saving against all-branches cover execution, with 52 percent maximum saving;
uniform model cells give 16 percent median and 59 percent maximum saving.
These are comparisons of validated envelopes, not new executions of commercial
defenses. The supplement retains actual alignment/all-branches/uniform/random
executions of the separate two-route microbenchmark.
Recompilation at output weights 0.25, 1, 2, 4, 8 and 16 accepts all 72 candidates
and leaves each of the twelve variants' shapes unchanged. This tests the given
profiles rather than establishing weight insensitivity for arbitrary programs.

For an observer restricted to exported reports, normalization achieves the same
fixed reporting abstraction with less execution work. Isolated prefill is useful
for consuming and auditing an envelope; it adds no input-length protection below
the wrapper. Historical latency-augmented attacks reach 100 percent balanced
accuracy on one shaped workflow. The present study therefore makes no timing
claim. All earlier application cohorts retain their original implementation
identities in the supplement and do not enlarge this homogeneous denominator.

\subsection{Concrete byte-bounded admission}\label{sec:concrete-admission}

The optional adapter in Section~\ref{sec:static-admission} is instantiated for
all twelve workflow variants. We scan every vocabulary entry of each pinned
GGUF with the complete native token-piece API: the two Llama vocabularies each
contain 128,256 entries and Qwen contains 151,936. All three have maximum raw
piece length 128 bytes. There are 231 pieces longer than 32 bytes in each Llama
vocabulary and 240 in Qwen. The former fixed-buffer wrapper returned an empty
piece for a 128-byte witness; the corrected decoder resizes the buffer and
retains the complete piece. No generic 32-byte bound is assumed.

The adapter explicitly bounds UTF-8 input buffers, visible responses and tool
results. A response buffer has size $128k$ for visible token cap $k$; it is
applied after decoding, because replacement decoding of invalid UTF-8 can expand
the byte count. Tool and input buffers use four times their declared token
annotation, as a chosen byte policy rather than a tokenizer conversion theorem.
Both original and shaped executions use the same buffers. The frontend derives
slot bounds from all source assignments and sizes each target with $E_v$.
It then exports the actual candidate and bounds to Lean's
\code{staticCertificate}, alongside independent source-bound validation.

For these byte-BPE tokenizers we declare $\alpha_m=1$ and $\delta_m=0$ on valid
UTF-8 complete requests, with automatic BOS insertion disabled. Every vocabulary
contains all 256 single-byte pieces. The native byte-splitting/merge algorithm
remains trusted~\cite{llamacpp2026vocab}; 2,058 adversarial and seeded complete
requests per model have no counterexample to the bound. Arbitrary isolated
non-UTF-8 byte strings are outside the source-string domain. A Lean lemma proves
that bounded token-piece lengths and a token cap bound the sum of raw piece
bytes; the adapter's post-decoding buffer enforces the actual string bound.
The scan and executable certificate do not verify native code or JSON parsing.

\begin{table}[t]
\centering\footnotesize
\caption{Actual source-derived admission candidates. All twelve pass core and
static admission checks at context 21,504. Guard counts refer to the full planned
input sets; maximum $T+N$ is taken over the emitted shape. Executed pairs select
the shortest and longest admitted UTF-8 input before inference.}
\label{tab:round2-admission}
\begin{tabular}{lrrr}\toprule
Variant & Guard admitted & Max. $T+N$ & Equal/shape pairs\\\midrule
\tabinput{generated/round2-admission.tex}
\bottomrule\end{tabular}
\end{table}

The guard admits \roundTwoGuard{} of the 1,684 planned inputs; two Calculator
inputs exceed the chosen byte policy. The 24 separate pairs complete with equal
full runtime stores and their certified shapes. Auditing reconciles
\roundTwoByteCalls{} local model calls, \roundTwoToolClips{} tool-buffer clips
and \roundTwoResponseClips{} response-buffer clips. Deliberately oversized inputs
are refused before any dispatch for all twelve variants. This is an implemented
bounded interface with explicit admission, not evidence that the unbounded
primary interface satisfies the same value invariant. The theorem covers all
guarded trajectories under its primitive contracts; finite executions test the
adapter without proving a physical failure refinement.
An earlier 32,768-context attempt completed 16 pairs before device-memory
allocation failed at the adaptive workflow. Those records remain separate:
token-context admission alone does not bound simultaneous device allocations.

\subsection{Single-batch native prefill}\label{sec:native-prefill}

We also implement a separate greedy prototype using one native
\code{llama\_decode} submission with 128 tokens. Useful request tokens occupy
sequence 0 and filler occupies sequence 1; each sequence has its own positions,
and only the last useful token requests logits. Subsequent decode submissions
contain one token. Sequence membership provides attention isolation through the
native batch interface~\cite{llamacpp2026batch}. This experiment runs on CPU with
two model families, three empty/short/multilingual prompt bodies, visible caps
4 and 16, and forced length 24. All attempts remain in Table~\ref{tab:round2-native}.

\begin{table}[t]
\centering\footnotesize
\caption{Native-prefill prototype, 12 attempted cases per configuration.
Single counts one exact 128-token prefill; Equal compares the visible output
with natural unpadded generation. Metadata counts cases in which sequence IDs
recover the useful request length. Cache removal discards filler sequence 1
after prefill.}
\label{tab:round2-native}
\begin{tabular}{lrrrr}\toprule
Native microbatch & Done & Single & Equal & Metadata\\\midrule
\tabinput{generated/round2-native.tex}
\bottomrule\end{tabular}
\end{table}

For each model, the natural native-size traces have four distinct values and
the shaped size projection has one: 128 followed by 23 single-token decode
submissions. Normalizing reports above this API cannot change those submissions.
Microbatch size 1 preserves all 12 visible outputs. At size 128, one Qwen case
differs. Removing the cover cache does not remove that mismatch. These
observations refute unqualified transparency for the latter configuration;
they do not prove transparency for every input at size 1.
Sequence and logit metadata still identify the useful-token subset in every
case. The prototype therefore supports a fixed native-size projection, while a
server observing full batch metadata retains request-length information.
No claim extends this projection to packets, cache traffic, elapsed time or
arbitrary physical failures.

\subsection{Categorical corrective-RAG ablation}

The primary Qwen grader mixes numeric-score instructions with a yes/no parser.
The earlier retained diagnosis found 484 zeros, 34 ones, five yes responses and
77 no responses among 600 grades; accepting numeric ones still sent all 150
questions to web search. We retain that primary port and evaluate a separate
contract: a categorical relevance prompt and a grammar permitting only
\code{\{"score":"yes"\}} or \code{\{"score":"no"\}}. The grammar is part
of the grader's model-contract identity. Answer generation, retrieval and the
rule ``search if any of four grades is not yes'' remain the same.

The ablation runs all 150 paired questions. All 600 original grades are valid
categorical JSON: \categoricalYes{} yes and \categoricalNo{} no. Web search is
selected for \categoricalSearch{} questions and omitted for
\categoricalNoSearch{}. Task accuracy is \categoricalTask{} percent under the
same inherited answer scorer: \categoricalCorrect{}/150 correct, compared with
\categoricalPrimaryCorrect{}/150 for the primary Qwen port. The paired comparison
has \categoricalImproved{} improved and \categoricalRegressed{} regressed answers.
The route distribution remains strongly imbalanced. All pairs preserve saved
fields, including the four grades, and have one shaped logical trace. Format compliance does not
establish grader correctness, and this combined prompt/grammar intervention
does not identify either component's individual causal contribution.
